\chapter{Hantavirus}

\section*{Definition and Overview}
Hantaviruses are a group of viruses carried by rodents that can cause serious disease in humans. They are spread through contact with rodent urine, droppings, or saliva, usually when contaminated dust is inhaled. There are two main forms of disease: hantavirus pulmonary syndrome (HPS), which causes severe lung failure, and haemorrhagic fever with renal syndrome (HFRS), which affects the kidneys. Hantavirus disease is rare in Australia, where the main risk is from rodents in rural and agricultural settings, but it is an important infection in parts of Asia, Europe, and the Americas. There is no specific treatment or vaccine, so prevention through rodent control and safe handling of rodent-contaminated areas is essential.

\section*{Epidemiology}
Hantavirus disease is rare in Australia. The Australian native hantavirus has not been shown to cause human disease, and the human-pathogenic forms are acquired mainly from imported rodents, particularly in port and shipping settings and in agricultural areas. Globally, thousands of cases occur each year: hantavirus pulmonary syndrome is concentrated in the Americas, and haemorrhagic fever with renal syndrome occurs mainly in Asia and Europe, with the highest numbers in China, Korea, and Russia. Risk is highest in rural and farming communities and in people who clean rodent-infested buildings.

\section*{Aetiology and Pathophysiology}
Hantaviruses infect humans through contact with infected rodents.

\begin{itemize}
    \item \textbf{The virus:} hantaviruses belong to the bunyavirus family, and different rodent species carry different strains
    \item \textbf{Transmission:} inhalation of dust contaminated with rodent urine, droppings, or saliva; direct contact with rodents; and rarely rodent bites
    \item \textbf{No person-to-person spread:} apart from a rare South American strain (Andes virus), humans do not pass it to each other
    \item \textbf{Incubation:} symptoms appear one to eight weeks after exposure
    \item \textbf{Hantavirus pulmonary syndrome:} the virus damages the lining of small blood vessels, causing fluid to leak into the lungs
    \item \textbf{Haemorrhagic fever with renal syndrome:} kidney inflammation and impaired function, with bleeding in severe cases
    \item \textbf{Pathology:} widespread inflammation of blood vessels (vasculitis) underlies both forms
\end{itemize}

\section*{Clinical Features}
The two forms have different patterns of illness.

\begin{itemize}
    \item \textbf{HPS early symptoms:} fever, muscle aches (especially in the thighs, hips, and back), headache, chills, nausea, and fatigue
    \item \textbf{HPS severe phase:} cough, breathlessness, and rapid progression to respiratory failure
    \item \textbf{HFRS:} sudden fever, headache, back pain, abdominal pain, nausea, and vomiting
    \item \textbf{HFRS kidney phase:} low urine output, bleeding, and in severe cases kidney failure
    \item \textbf{Recovery:} gradual over weeks to months
    \item \textbf{Exposure history:} contact with rodents or rodent-contaminated environments
\end{itemize}

\section*{Differential Diagnosis}
The early symptoms of hantavirus resemble many other infections.

\begin{itemize}
    \item Influenza and other respiratory viruses
    \item COVID-19
    \item Leptospirosis, another rodent-associated infection
    \item Q fever
    \item Other causes of pneumonia, including bacterial and viral pneumonia
    \item Sepsis from other causes
    \item Other causes of kidney failure and bleeding
    \item Dengue and other mosquito-borne infections
\end{itemize}

\section*{When to Seek Medical Care}
Seek medical care promptly if you develop fever, muscle aches, and breathlessness after exposure to rodents or rodent-contaminated environments, such as cleaning a shed, barn, or infested building.

\textbf{Call 000 (or local emergency services) for:}
\begin{itemize}
    \item Difficulty breathing or rapidly worsening breathlessness
    \item Coughing up blood
    \item Very low urine output or signs of kidney failure
    \item Confusion, seizures, or collapse
    \item Any severe, rapidly progressive illness
\end{itemize}

\section*{Diagnosis and Investigations}
Hantavirus is diagnosed by blood tests and clinical assessment.

\begin{itemize}
    \item \textbf{Blood tests:} antibodies to the virus appear during the illness
    \item \textbf{PCR testing:} detects the virus's genetic material in the blood early in the illness
    \item \textbf{Full blood count:} low platelets and other changes are characteristic
    \item \textbf{Kidney function tests:} for HFRS
    \item \textbf{Chest X-ray and oxygen monitoring:} for HPS
    \item \textbf{Exposure history:} crucial for considering the diagnosis
    \item \textbf{Notification:} hantavirus infection is notifiable in Australia
\end{itemize}

\section*{Management}
There is no specific antiviral treatment, so care is supportive and intensive.

\begin{itemize}
    \item \textbf{Hospital care:} both forms of hantavirus disease require hospital admission
    \item \textbf{HPS:} intensive care with oxygen and, in severe cases, mechanical ventilation
    \item \textbf{HFRS:} kidney support, including dialysis in severe cases
    \item \textbf{Fluids and electrolytes:} careful management
    \item \textbf{Bleeding management:} transfusions where needed
    \item \textbf{Pain and fever relief:} supportive medicines
    \item \textbf{Recovery support:} gradual return to activity over weeks to months
    \item \textbf{Prevention of spread:} standard infection control, as person-to-person spread is very rare
\end{itemize}

\section*{Complications}
\begin{itemize}
    \item Respiratory failure requiring ventilation in HPS
    \item Kidney failure requiring dialysis in HFRS
    \item Bleeding complications
    \item Shock
    \item Death: HPS is fatal in around 30--40% of cases, and HFRS in 1--5%
    \item Prolonged fatigue and reduced kidney function in survivors
    \item Complications of intensive care
\end{itemize}

\section*{Prognosis and Outcomes}
The outcome depends on the form of disease and the speed of treatment. Hantavirus pulmonary syndrome is serious, with a fatality rate of around 30--40%, though intensive care has improved survival. Haemorrhagic fever with renal syndrome is fatal in a small percentage of cases, and most people recover, though kidney function may take time to return to normal. Because the disease is rare in Australia and there is no treatment or vaccine, the emphasis is on prevention through rodent control and safe practices when working in rodent-infested areas.

\section*{Prevention and Self-Care}
\begin{itemize}
    \item Control rodents around homes, farms, and workplaces
    \item Seal holes and gaps that allow rodents to enter buildings
    \item Store food, grain, and animal feed in rodent-proof containers
    \item When cleaning rodent-infested areas, wear gloves and a mask
    \item Dampen dusty areas with water before sweeping to avoid inhaling contaminated dust
    \item Disinfect surfaces contaminated by rodents with bleach
    \item Remove rodent carcasses and droppings safely, using gloves and a mask
    \item Avoid sleeping in areas where rodents are active
    \item Seek medical care promptly if you develop fever and breathlessness after rodent exposure
\end{itemize}

\section*{Sources and Further Reading}
The following reputable sources provide further information for healthcare students and patients seeking reliable, current advice about hantavirus.

\begin{itemize}
    \item Australian Government Department of Health and Aged Care. \textit{Hantavirus fact sheet.} https://www.health.gov.au/
    \item Healthdirect Australia. \textit{Hantavirus.} https://www.healthdirect.gov.au/
    \item NSW Health. \textit{Hantavirus fact sheet.} https://www.health.nsw.gov.au/
    \item Centers for Disease Control and Prevention. \textit{Hantavirus.} https://www.cdc.gov/
    \item World Health Organization. \textit{Hantavirus.} https://www.who.int/
    \item European Centre for Disease Prevention and Control. \textit{Hantavirus infections.} https://www.ecdc.europa.eu/
\end{itemize}
